\begin{lexlean}{Tapwire.Invariants}
\useglossary{lexlean.std.bool@1.2.0}
\useglossary{lexlean.std.nat@1.2.0}
\useglossary{tapwire.titles@1.0.0}
\importmodule{Tapwire.Capture}
\importmodule{Tapwire.PcapNg}
\importmodule{Tapwire.Dissector}
\importmodule{Tapwire.Filter}
\title{Tapwire formal invariants}
\begin{semanticmodule}
theorem{
  name: capture_limits_bound_packet_length,
  parameters: [],
  statement: (((project (Tapwire.Capture::defaultCaptureLimits) maxPacketLength) <=? (project (Tapwire.Capture::defaultCaptureLimits) snaplen)) = true),
  proof: reflexivity{}
}

theorem{
  name: pcapng_block_alignment_is_enforced,
  parameters: [],
  statement: (((Tapwire.PcapNg::isAlignedTo32BitBoundary (Tapwire.PcapNg::alignedEnhancedPacketBlock)) && (not (Tapwire.PcapNg::isAlignedTo32BitBoundary (Tapwire.PcapNg::misalignedEnhancedPacketBlock)))) = true),
  proof: reflexivity{}
}

theorem{
  name: dissector_fields_respect_parent_and_depth,
  parameters: [],
  statement: (((Tapwire.Dissector::isValidFieldInstance (Tapwire.Dissector::ipSourceField)) && (Tapwire.Dissector::isValidFieldInstance (Tapwire.Dissector::tcpPortField))) = true),
  proof: reflexivity{}
}

theorem{
  name: display_filter_is_bounded,
  parameters: [],
  statement: (((Tapwire.Filter::filterNodeCount (Tapwire.Filter::tcpPortFilter)) <=? 64) = true),
  proof: reflexivity{}
}
\end{semanticmodule}
\end{lexlean}
